\subsection{Sample size optimization}
\label{subsec:sample_optimization}

For a fixed chamber assignment \(\mathbf{x}\), this stage changes the
sample vector \(\mathbf{s}\).
An extra sample adds one pull-down operation but can let later tests
of that product run in parallel.
Removing a sample does the opposite: pull-down work falls, and so
does the chance of overlap.
Which direction helps is not known in advance, so both are tried
under the same assignment.

The search never uses fewer than three samples per product.
A single physical unit would place every test of that product on the
same sample, so those tests cannot overlap, and a one-unit plan
would give no replication across units.
The laboratory therefore treats several units as the minimum for a
valid test plan.
Three is the industrial baseline: the local search starts there and
does not go below it.
An upper bound \(s_{\max}\) only keeps the plan from using an
unrealistic number of prototypes.

The neighborhood changes one product at a time by one or two
samples, provided the new count stays between \(s_{\min}\) and
\(s_{\max}\).
For each product the improving trial with the best tardiness---or,
when tardiness is equal, the smaller sample count---is kept.
After this single-product pass, a Late--Early Rebalance pairs late
products, ordered by decreasing tardiness, with early products,
ordered by decreasing slack, up to \(K_{\mathrm{reb}}\) pairs.
Each pair moves one sample from an early product to a late product
without leaving the bounds.
The paired move is accepted if total tardiness does not increase;
when tardiness is unchanged, the total number of samples must not
increase.
Every trial is decoded by the same earliest-due-date (EDD)
scheduler used in the outer framework.
Algorithm~\ref{alg:sample_ls} states the procedure.

\begin{algorithm}[H]
\caption{Sample Size Optimization}
\label{alg:sample_ls}
\begin{algorithmic}[1]
\State \textbf{Input:} chamber assignment \(\mathbf{x}\), sample vector \(\mathbf{s}\)
\State \textbf{Output:} \((\mathbf{s},\Pi,T^{tot})\)
\State Restrict every \(s_p\) to \([s_{\min},s_{\max}]\) \label{step:ls-init}
\State \((\Pi,T^{tot})\leftarrow\text{Schedule}(\mathbf{x},\mathbf{s})\) \label{step:ls-eval}
\For{\(k\leftarrow 1\) to \(I_{\max}^{(2)}\)} \label{step:ls-loop}
    \For{each product \(p\in\mathcal{P}\)}
        \State Try changing \(s_p\) by one or two samples within the bounds \label{step:ls-nbhd}
        \If{the best of those trials improves the incumbent} \label{step:ls-accept}
            \State Update \(s_p\) and set \((\Pi,T^{tot})\leftarrow\text{Schedule}(\mathbf{x},\mathbf{s})\)
        \EndIf
    \EndFor
    \State \(\mathbf{s}_{reb}\leftarrow\text{LateEarlyRebalance}(\mathbf{s},K_{\mathrm{reb}})\) \label{step:ls-reb}
    \State \((\Pi_{reb},T^{tot}_{reb})\leftarrow\text{Schedule}(\mathbf{x},\mathbf{s}_{reb})\)
    \If{the rebalance does not worsen the incumbent}
        \State \((\mathbf{s},\Pi,T^{tot})\leftarrow(\mathbf{s}_{reb},\Pi_{reb},T^{tot}_{reb})\)
    \EndIf
    \If{no move was accepted in this pass}
        \State \Return \((\mathbf{s},\Pi,T^{tot})\) \label{step:ls-stop}
    \EndIf
\EndFor
\State \Return \((\mathbf{s},\Pi,T^{tot})\)
\end{algorithmic}
\end{algorithm}

Step~\ref{step:ls-init} enforces the three-sample floor and the
upper bound.
Step~\ref{step:ls-eval} builds the EDD schedule of the current
vector; this is the incumbent of the local search.
Each pass (Step~\ref{step:ls-loop}) first tries the one-product
moves (Step~\ref{step:ls-nbhd}) and accepts a product's best
improving trial (Step~\ref{step:ls-accept}).
Step~\ref{step:ls-reb} then shifts one sample from an early product
to a late product.
That move need only not worsen tardiness; the single-product moves
must strictly improve the incumbent under the same rule as the outer
framework.
The search stops when a pass accepts nothing
(Step~\ref{step:ls-stop}) or when the pass limit \(I_{\max}^{(2)}\)
is reached.
